% Detailed publication captions for the five validation panels of
% the Epistemological Mode-Methodology Equivalence manuscript.
% Each caption describes axes, ranges, scales, the theoretical formula,
% the relevant theorem reference, and the measured precision against
% the closed-form prediction.

\begin{figure*}[!htbp]
    \centering
    \includegraphics[width=\textwidth]{validation/figures/panel_1.png}
    \caption{\textbf{Receiver Floor and Cell-Truth.}
    (\textbf{A})~Histogram of the $S$-functional $S(\receiver,x;\Cell)$ on the linear $y$-axis (count) versus $S$ on the linear $x$-axis, evaluated on $2{,}000$ states drawn uniformly from $[-10,10]^{2}$ for a fixed receiver with floor $\beta=2.5$ and an action-cell of tolerance $\tau=1.0$ centred at the origin. The empirical minimum of $S$ across all $2{,}000$ states equals $\beta=2.5$ (crimson dashed reference), confirming the floor bound $S\ge\beta>0$ from Theorem~\ref{thm:floor-receiver} (Floor Theorem for Receivers).
    (\textbf{B})~Cell-truth verification: $S$ on the linear $y$-axis ($\approx 1.5$ to $\approx 8$) versus $|x|$ on the linear $x$-axis ($0$ to $8$), separating $200$ states drawn from inside the cell (seagreen, radii $\le 1.8$) and $200$ from outside (orange, radii $\ge 2.5$). All inside-cell points collapse onto the floor line $S=\beta=1.5$ (crimson dashed), independently of $|x|$, while outside-cell points lie strictly above by the cell-distance $d(x,\Cell)$. This is the empirical content of Theorem~\ref{thm:cell-truth} (Cell-Truth): all states inside an action-equivalence cell are $S$-indistinguishable.
    (\textbf{C})~Floor attainment as a function of receiver floor: measured $S(\receiver,x;\Cell)$ at a fixed in-cell state $x=(0.5,0)$ on the linear $y$-axis ($0$ to $5$) versus $\beta$ on the linear $x$-axis ($0.1$ to $5$), $30$ values of $\beta$. Measured (steelblue circles) coincides exactly with the theoretical line $S=\beta$ (crimson, slope~$1$). Verifies the attainment clause of Theorem~\ref{thm:floor-receiver} across the full $\beta$-range; experiment E03 reports relative error $0.00$.
    (\textbf{D})~Three-dimensional surface of $S(\receiver,x;\Cell)$ over the state slice $(x,y)\in[-3,3]^{2}$, $25\times 25$ grid, with three superimposed surfaces for $\beta\in\{0.5,1.5,2.5\}$. Each surface is flat at $S=\beta$ over the cell interior (the disc of radius $1$ around the origin) and rises linearly with cell distance $d(x,\Cell)$ outside it; the three surfaces are vertical translates separated by the floor difference $\Delta\beta=1.0$, demonstrating that the cell shape is invariant of $\beta$ while the elevation tracks the floor.}
    \label{fig:panel1-receiver-floor}
\end{figure*}

\begin{figure*}[!htbp]
    \centering
    \includegraphics[width=\textwidth]{validation/figures/panel_2.png}
    \caption{\textbf{Representational Invariance: Oscillatory, Categorical, Partition Encodings.}
    (\textbf{A})~Oscillatory invariance: $S$ after applying an isometric oscillatory encoding $\phi_{o}$ (rotation by $0.7$~rad) on the linear $y$-axis versus $S$ in the original representation on the linear $x$-axis, $100$ random states from $[-3,3]^{2}$ with action-cell of tolerance $0.7$ centred at $(1,0.5)$, receiver floor $\beta=1.0$. Steelblue points fall exactly on the diagonal $y=x$ (crimson dashed), confirming that the oscillatory encoding preserves $S$ to machine precision; experiment E07 reports relative error $0.00$. Empirical content of Theorem~\ref{thm:rep-inv} (Representational Invariance), oscillatory clause.
    (\textbf{B})~Categorical invariance: same protocol with the categorical encoding $\phi_{c}$ (signed coordinate permutation in $\mathbb{R}^{3}$). Seagreen points coincide with $y=x$, max absolute deviation $0.00$; experiment E08 reports relative error $0.00$. Empirical content of Theorem~\ref{thm:rep-inv}, categorical clause.
    (\textbf{C})~Partition invariance: same protocol with the partition encoding $\phi_{p}$ (sign flip). Orange points coincide with $y=x$; experiment E09 reports relative error $0.00$. Empirical content of Theorem~\ref{thm:rep-inv}, partition clause.
    (\textbf{D})~Three-dimensional triple-equivalence manifold: $S$ in the categorical encoding on the $z$-axis versus $S$ in the oscillatory encoding on the $y$-axis versus $S$ in the original encoding on the $x$-axis, $100$ scatter points (steelblue), with the triple-equivalence diagonal $\{(t,t,t):t\in[0,5]\}$ overlaid (crimson dashed). All points lie on the diagonal, demonstrating that the three encodings yield the same $S$ value simultaneously --- the geometric realization of Corollary~\ref{cor:triple} (Triple Equivalence at the level of $S$): choice of representation is computational, not epistemic.}
    \label{fig:panel2-rep-invariance}
\end{figure*}

\begin{figure*}[!htbp]
    \centering
    \includegraphics[width=\textwidth]{validation/figures/panel_3.png}
    \caption{\textbf{Layered Receivers and Mode Non-Privilege.}
    (\textbf{A})~Per-layer floors and aggregate floor: bar chart of layer floors $\beta_{i}$ on the linear $y$-axis ($0$ to $6$) for five layers $L_{1},\dots,L_{5}$ with floors $(0.3,0.8,1.5,3.0,6.0)$. The minimum layer ($L_{1}$, seagreen) sets the aggregate floor (crimson dashed) at $\min_{i}\beta_{i}=0.3$. Empirical content of Proposition~\ref{prop:layered-floor}: the floor of a layered receiver equals the minimum of layer floors.
    (\textbf{B})~$S$-contour map for a two-layer receiver: filled contour of $S(\receiver_{\bullet},x;\Cell)$ over the state plane $(x,y)\in[-3,3]^{2}$, $60\times 60$ grid, viridis colourmap, with pre-decoder floor $\beta_{1}=0.4$ and decoder floor $\beta_{2}=2.5$. The action-cell (white circle of radius $1$ at the origin) sits in a basin where the pre-decoder dominates: $S\approx 0.4$ throughout the firing region. The decoder floor would have $S\ge 2.5>\tau(\Cell)=1.0$, so the decoder alone could not reach the cell. Empirical content of Theorem~\ref{thm:mode-nonpriv} (Mode Non-Privilege): pre-decoder layer dominates when $\beta_{1}<\tau(\Cell)<\beta_{2}$.
    (\textbf{C})~Reachability: fraction of states (out of $500$ uniform samples from $[-3,3]^{2}$) for which a single layer reaches the action-cell, on the linear $y$-axis ($0$ to $1$), versus reflex floor on the linear $x$-axis ($\{0.2,0.5,1.0,2.0\}$). Reflex layer reach (seagreen circles) drops as the reflex floor approaches the cell tolerance $\tau=1.0$; decoder layer reach (orange squares, fixed $\beta_{2}=2.5>\tau$) is identically zero across all reflex-floor values. Demonstrates Corollary~\ref{cor:not-unique} (Knowledge is not the unique mode): the action-cell is reachable through reflex even when decoder reach is zero.
    (\textbf{D})~Three-dimensional surface of the aggregate floor $\beta_{\rm agg}=\min(\beta_{\rm reflex},\beta_{\rm decoder})$ over $(\beta_{\rm reflex},\beta_{\rm decoder})\in[0.1,3.0]\times[0.1,5.0]$, $20\times 20$ grid, viridis colourmap. The surface exhibits the characteristic min-aggregation ridge along the diagonal $\beta_{\rm reflex}=\beta_{\rm decoder}$: for $\beta_{\rm reflex}<\beta_{\rm decoder}$ the surface is independent of $\beta_{\rm decoder}$ (reflex-dominated regime); for $\beta_{\rm decoder}<\beta_{\rm reflex}$ the surface is independent of $\beta_{\rm reflex}$ (decoder-dominated regime). The minimum-of-two structure of Proposition~\ref{prop:layered-floor}.}
    \label{fig:panel3-layered-receivers}
\end{figure*}

\begin{figure*}[!htbp]
    \centering
    \includegraphics[width=\textwidth]{validation/figures/panel_4.png}
    \caption{\textbf{Methodological Floor and the Catalytic Fixed Point.}
    (\textbf{A})~$S$-trajectory across $50$ iterations of three methodologies with common dispersion $\sigma=0.5$ and contraction $\kappa\in\{0.3,0.5,0.7\}$, starting from $s_{0}=50$. $S$ on the log $y$-axis ($\sim 10^{-1}$ to $10^{2}$) versus iteration count on the linear $x$-axis ($0$ to $50$). Each trajectory contracts geometrically and asymptotes to its predicted methodological floor $\Sfloor(\Method)=\sigma\kappa/(1-\kappa)\in\{0.214,0.500,1.167\}$ for $\kappa\in\{0.3,0.5,0.7\}$. Empirical content of Theorem~\ref{thm:method-floor} (Methodological Floor): no number of iterations of $\Method$ alone reduces $S$ below $\Sfloor(\Method)$; convergence to floor verified to relative error $\le 1.85\times 10^{-16}$ in experiment E16.
    (\textbf{B})~Methodological floor $\Sfloor(\Method)$ on the linear $y$-axis versus contraction $\kappa$ on the linear $x$-axis ($0.05$ to $0.95$), $30$ values, fixed dispersion $\sigma=0.5$. Steelblue curve traces the closed form $\Sfloor=\sigma\kappa/(1-\kappa)$ from Definition~\ref{def:method}. The floor diverges as $\kappa\to 1^{-}$, encoding that highly-correlated methodologies (low contraction, large $\kappa$) carry large irreducible floors.
    (\textbf{C})~Methodological floor $\Sfloor(\Method)$ on the linear $y$-axis versus dispersion $\sigma$ on the linear $x$-axis ($0$ to $2$), $30$ values, three curves for $\kappa\in\{0.3,0.5,0.7\}$. All curves are linear in $\sigma$ with slope $\kappa/(1-\kappa)\in\{0.428,1.000,2.333\}$, so the floor scales linearly with per-iteration dispersion. Empirical evidence for the closed form $\Sfloor=\sigma\kappa/(1-\kappa)$.
    (\textbf{D})~Three-dimensional surface of $\Sfloor(\Method)$ over the parameter region $(\kappa,\sigma)\in[0.05,0.9]\times[0.05,1.5]$, $30\times 30$ grid, viridis colourmap. The surface is monotonically increasing in both arguments, vanishes along the lower-left edge ($\kappa\sigma\to 0$), and rises sharply along the right edge ($\kappa\to 1^{-}$). This is the joint characterization of Theorem~\ref{thm:method-floor}: methodologies are characterized by a two-parameter floor that no methodology with $\sigma>0,\kappa>0$ can avoid.}
    \label{fig:panel4-methodological-floor}
\end{figure*}

\begin{figure*}[!htbp]
    \centering
    \includegraphics[width=\textwidth]{validation/figures/panel_5.png}
    \caption{\textbf{Mode--Methodology Equivalence and the Distribution Theorem.}
    (\textbf{A})~Composite floor: measured $\Sfloor(\receiver\circ\Method)$ on the linear $y$-axis versus predicted $\Sfloor=\Sfloor(\receiver)+\Sfloor(\Method)-\Sfloor(\receiver)\Sfloor(\Method)/\Sigma$ on the linear $x$-axis, $40$ random pairs $(\Sfloor(\receiver),\Sfloor(\Method))$ drawn uniformly from $[0.1,5]^{2}$, normalization $\Sigma=100$. Steelblue scatter falls exactly on the diagonal $y=x$ (crimson dashed), confirming the closed-form composition law of Theorem~\ref{thm:mode-meth-equiv} (Mode--Methodology Equivalence). Maximum relative error across the $40$ pairs: $9.37\times 10^{-15}$, IEEE-754 rounding noise (experiments E20--E22).
    (\textbf{B})~Multiplicative reduction of composite floor: composite $\Sfloor$ on the linear $y$-axis ($0$ to $\sim 10$) versus stack length $n$ on the linear $x$-axis ($1$ to $10$), four curves for per-stack floor $\beta_{\star}\in\{0.5,1.0,1.5,2.0\}$. Each curve traces $\Sigma(1-(1-\beta_{\star}/\Sigma)^{n})$ from Corollary~\ref{cor:stack}; composite floor grows monotonically with $n$ but the marginal contribution decays geometrically. Verifies Theorem~\ref{thm:catcomp} (Catalytic composition law) for repeated stacking.
    (\textbf{C})~Anti-monopoly: maximum knowledge $\sup_{x\in X_{0}}\mathrm{Know}_{\receiver}(x)=\Sigma-\Sfloor(\receiver)$ on the linear $y$-axis (numerator $\Sigma=100$) versus receiver floor on the linear $x$-axis ($0.1$ to $10$), $30$ values. Steelblue line traces the upper bound $\Sigma-\beta$; orange circles report the empirical maximum knowledge over $100$ random states with $\eps$-ball cells. Measured maximum coincides with bound at every $\beta$: experiment E24 reports relative error $0.00$. Empirical content of Theorem~\ref{thm:distribution} (Distribution): no bounded receiver can hold full knowledge ($\Sigma$) of any subject; the deficit equals the receiver's floor.
    (\textbf{D})~Three-dimensional composite-floor surface $\Sfloor(\receiver\circ\Method)=\Sfloor(\receiver)+\Sfloor(\Method)-\Sfloor(\receiver)\Sfloor(\Method)/\Sigma$ over $(\Sfloor(\receiver),\Sfloor(\Method))\in[0.1,10]^{2}$, $25\times 25$ grid, viridis colourmap. The surface is symmetric in $\Sfloor(\receiver)\leftrightarrow\Sfloor(\Method)$ (line of constant elevation along the diagonal), monotone in each coordinate, and slightly sub-additive due to the $-\Sfloor(\receiver)\Sfloor(\Method)/\Sigma$ correction. The symmetry is the geometric expression of Remark~\ref{rem:symmetry}: at the level of attainable $S$, swapping a receiver-layer for a methodology with the same floor leaves the composite floor invariant.}
    \label{fig:panel5-mode-methodology}
\end{figure*}
